% ============================================================
\subsection{Feedback Loop and Adaptation}
\label{subsec:feedback-loop}
% ============================================================

The components described so far are sufficient to serve a single ad on a
single turn, but they are not sufficient to make the system \emph{learn} from
its own behaviour. This subsection describes the asynchronous feedback loop
that observes user reactions to served ads, distinguishes genuine engagement
from noise, and feeds that signal back into the matching pipeline. We regard
this loop as one of the central contributions of the project, because it is
the part of the design that takes seriously the observation that
conversational ads cannot be evaluated by clicks alone.

\subsubsection{Click Tracking, Event Logging, and Budget Accounting}
\label{subsubsec:click-tracking}

Every ad that appears in a user-facing response carries a call-to-action URL.
Rather than expose the advertiser's destination URL directly, the orchestrator
rewrites it at synthesis time into a tracked redirect of the form
\texttt{/track/click?ad\_id=X\&session\_id=Y\&redirect=Z}, where $Z$ is the
URL-encoded original destination. When the user clicks the link, the
\texttt{/track/click} endpoint writes a \texttt{click} event to the
\texttt{events} table and issues a 302 redirect to $Z$. The user perceives a
single click; the system gains a logged interaction tied to the originating
session and ad.

The event model is uniform across the four event types we record:
\texttt{impression}, \texttt{click}, \texttt{engagement}, and
\texttt{conversion}. Each row carries an \texttt{event\_id}, an
\texttt{ad\_id}, a \texttt{session\_id}, a \texttt{turn\_index}, a UTC
timestamp, and a JSONB \texttt{payload} column for type-specific fields
(e.g.\ the classifier confidence for an engagement event, or the referring
turn for a click). Storing the type-specific fields in JSONB rather than in
separate tables keeps the analytics queries simple --- a single
\texttt{GROUP BY event\_type} suffices to compute most dashboard metrics ---
at the cost of giving up some compile-time schema enforcement, which we
judge an acceptable trade for a research prototype.

Budget accounting is tied to the impression event. When the orchestrator
logs an impression, it does so in the same PostgreSQL transaction that
decrements \texttt{remaining\_budget} on the corresponding ad row by
$\texttt{bid\_cpm}/1000$. Because the event insert and the budget decrement
are atomic, the system maintains the invariant that the sum of charged
impressions per ad, multiplied by its CPM, equals the total amount drawn
down from its budget. This is the property that lets the analytics layer
report revenue numbers that reconcile with the database rather than drift
from it over time.

\subsubsection{LLM-Judge-Based Engagement Detection}
\label{subsubsec:llm-judge-engagement}

Clicks are a poor primary signal in a conversational setting. A user who is
intrigued by a sponsored mention is far more likely to ask a follow-up
question about the product than to click an inline link, and yet that
follow-up is exactly the kind of behaviour that classical ad analytics is
blind to. The engagement detection module is the system's answer to this
mismatch: it inspects the user's \emph{next} turn after an ad was served and
decides whether that turn constitutes engagement with the advertised product.

We implemented this in two stages. The first stage is a cheap keyword and
entity matcher: if the next user message contains the ad's product name,
brand, or any of the entities extracted from its title, we tentatively flag
the turn as engaged. This catches the easy cases (``tell me more about the
Pegasus 41'') at essentially zero cost. The second stage, invoked only when
the first stage is ambiguous, is a lightweight LLM judge prompted with the served ad,
the previous assistant turn, and the new user turn, and asked to return a
boolean \texttt{engaged} field together with a short rationale and a
confidence score in $[0,1]$.

The two-stage design is a deliberate latency and cost optimisation. The
keyword stage handles the majority of true positives without any model call;
the LLM stage handles the genuinely ambiguous cases, where the user
paraphrases (``how durable are those shoes you mentioned?'') or asks an
implicit follow-up. When the LLM judge returns \texttt{engaged = true} with
confidence above a threshold $\tau = 0.7$, an \texttt{engagement} event is
written to the events table with the rationale stored in the JSONB payload.
We treat this event as a strictly stronger signal than a click, both in the
analytics dashboard and, as described next, in the adaptive layer that
modulates future ad selection.

It is worth being explicit about what the engagement detector is \emph{not}.
It is not a sentiment classifier, and it does not try to decide whether the
user liked the ad. It only decides whether the user's next turn is
\emph{about} the advertised product. Whether that engagement is positive,
negative, or neutral is left to downstream analysis; what matters for the
feedback loop is that the user's attention measurably moved towards the ad.

\subsubsection{Adaptive Ad Context Optimisation}
\label{subsubsec:adaptive-optimisation}

The final piece of the feedback loop closes it: engagement events are not
merely logged for reporting, they actively reshape what the matching
pipeline does on subsequent turns. We call this \emph{adaptive ad context
optimisation}, and it operates at two granularities.

\paragraph{Within-session adaptation.} When an engagement event is recorded
for an ad $a$ in session $s$, the orchestrator updates the Redis session
state for $s$ in two ways. First, the topics and entities associated with
$a$ are merged into the session's \texttt{accumulated\_topics} list, which
biases the next round of context extraction towards the same product
neighbourhood. Second, the session's \texttt{purchase\_signal} is bumped
upward by a fixed increment $\Delta = 0.15$, capped at $1.0$. The effect on
the very next turn is concrete and measurable: the query embedding shifts
toward the engaged-with topic, and the re-ranker's purchase-signal term
increases, jointly making it more likely that a related and complementary
ad surfaces if one exists in the inventory.

\paragraph{Cross-session adaptation.} At the inventory level, engagement and
click events feed a slow-moving per-ad quality score
$q_a \in [0,1]$, computed as a smoothed engagement-plus-click rate over the
ad's recent impressions. This score is not part of the per-request scoring
formula in Section~\ref{subsubsec:reranking-business-logic} --- introducing
it there would risk a runaway feedback loop in which a few early lucky
impressions monopolise future serving --- but it is exposed to the analytics
dashboard and used to surface under-performing ads to the (currently
manual) inventory review process. A learned re-ranker that incorporates
$q_a$ directly, with appropriate exploration, is the natural next step and
is sketched in Section~\ref{subsec:future-work}.

Taken together, these two granularities give the system a feedback loop
with two distinct time constants: a fast, within-session loop that adapts
the next turn's ad selection to what the user just engaged with, and a
slow, cross-session loop that lets the inventory as a whole drift toward
ads that consistently earn user attention. Neither loop requires
labelled training data or a separate model training pipeline; both are
driven entirely by the events that the system already records as part of
its normal operation.